\documentclass[11pt]{article}
\usepackage{amsmath,amssymb,amsthm}
\newtheorem{proposition}{Proposition}
\begin{document}
\title{An auxiliary butterfly switch repairs a strict positive obstruction}
\author{}\date{}\maketitle
\paragraph{Status.}
This is an exact finite-dimensional construction and an all-small-scale
continuation certified by rational arithmetic. It is not Lean formalized,
not an asymptotic compiler, and not a universality theorem for underloaded
butterflies. It shows that freeing a routing switch can be strictly
stronger than deleting its sign constraint at the original routing.

Use the binary butterfly with $H=16$ wires and switches $(k,s,g)$ as in
\texttt{butterfly-rectangle-decoder.tex}, ordered by increasing stage $k$,
then the order in its companion checker. Let
\[
 p_y(\epsilon)=\sum_{j=0}^3 y_j\epsilon^{2^j},\qquad
 t_{k,s}=1-p_s(\epsilon).
\]
For a switch word $\theta$, let $\pi_\theta$ be its routed permutation and
\[
 u_{k,g}(h)=\begin{cases}
  1,&h\mathbin{\mathrm{div}}2^{k+1}=g,\ h_k=0,\\
 -1,&h\mathbin{\mathrm{div}}2^{k+1}=g,\ h_k=1,\\
 0,&\text{otherwise}.
 \end{cases}
\]
For residues $c_y$, the signed query at switch $e=(k,s,g)$ is
\[
 S_e(\theta,c)=
 (1-2\theta_e)\sum_{y=0}^{15}
 \frac{u_{k,g}(\pi_\theta^{-1}(y))c_y}{1-p_s-p_y}.
\]
The following statement uses no block biases.

\begin{proposition}
Let the four consecutive stage words be
\[
 \theta=(10001101\ ;\ 11011111\ ;\ 00001111\ ;\ 00011110).
\]
Omit the query at the auxiliary switch $e_*=(3,3,0)$, which has index
$27$ when indices start at zero. For every sufficiently small positive
$\epsilon$, there is no nonnegative residue vector making all remaining
$31$ scores $S_e(\theta,c)$ strictly positive.

Let $\widehat\theta$ differ only at $e_*$, changing that unused routing
bit from $1$ to $0$. For every sufficiently small positive $\epsilon$,
there is a strictly positive residue vector $c(\epsilon)$ making all
$31$ retained scores $S_e(\widehat\theta,c(\epsilon))$ strictly positive.
The labels at all retained queries are unchanged.
\end{proposition}
\begin{proof}
The first assertion is already certified by
\texttt{butterfly-adaptive-haar-positive-h16-dual.json} and
\texttt{butterfly-adaptive-haar-positive-dual-check.py}.
Its nonnegative dependence uses query indices
\[
 5,10,12,17,18,19,20,21,23,24,25,26,28,29,30,31
\]
and the nonnegative-residue row for coordinate $1$. In particular it
omits query $27$. Its exact positive leading Laurent coefficients prove
infeasibility on an entire punctured interval, as established in the
existing obstruction writeup.

For the second assertion, form the $31\times16$ rational matrix $A$
whose rows are the retained signed queries for $\widehat\theta$.
Select its rows with original query indices
\[
 5,14,15,17,18,19,20,21,22,24,25,26,28,29,30,31
\]
to form a square matrix $B(\epsilon)$. Define
$c(\epsilon)=B(\epsilon)^{-1}\boldsymbol1$.
The saved checker verifies an exact rational witness at
$\epsilon=1/100$, and separately computes the leading Laurent terms
of this rational-function solution. The elimination pivot orders are
\[
 0,0,0,0,0,0,2,2,2,6,6,6,8,18,20,25.
\]
In particular the determinant is not identically zero. Every one of
the $31$ coordinates of $A c$ and the $16$ coordinates of $c$ has a
strictly positive leading Laurent coefficient. Therefore they are all
strictly positive for sufficiently small positive $\epsilon$.

The companion JSON records the exact finite witness, the chosen basis,
and every leading coefficient and exponent. The checker reconstructs
all matrix entries independently, verifies the finite inequalities,
and recomputes the rational-function leading coefficients using guarded
exact power-series arithmetic at precision $512$. Only pivots with a
known nonzero leading term are inverted. Unresolved series entries
are required to have sufficient remaining precision before choosing a
pivot, and the final coefficients and matrix residuals are checked.
Thus the continuation does not follow from numerical sampling.
\end{proof}

\paragraph{Remaining question.}
This example retains $31$ of $32$ switch queries, but concerns one
particular word. A growing construction would need a fixed selected
set of order $H\log H$ switches such that every assignment of labels
on that set admits auxiliary switch completions and suitable residues.
The example establishes neither that property nor a positive-density
version. It motivates testing the existential completion formulation
instead of requiring every full routing word to work.
\end{document}
